% Generated draft: requires mathematical and source review.
\documentclass[11pt]{article}
\usepackage[margin=25mm]{geometry}
\usepackage{fontspec}
\setmainfont{Latin Modern Roman}
\usepackage{amsmath,amssymb,mathtools,unicode-math}
\setmathfont{Latin Modern Math}
\usepackage{enumitem,needspace,xurl,hyperref}
\setlength{\parindent}{0pt}
\setlength{\parskip}{5pt}
\setlength{\emergencystretch}{3em}
\newcommand{\N}{\mathbb{N}}
\newcommand{\Z}{\mathbb{Z}}
\newcommand{\Q}{\mathbb{Q}}
\newcommand{\R}{\mathbb{R}}
\newcommand{\C}{\mathbb{C}}
\newcommand{\F}{\mathbb{F}}
\newcommand{\A}{\mathbb{A}}
\newcommand{\PP}{\mathbb{P}}
\newcommand{\E}{\mathbb{E}}
\newcommand{\eps}{\varepsilon}
\newcommand{\dd}{\,\mathrm{d}}
\begin{document}
\section*{UnsolvedMath ID: 30006290}
\textbf{Source problem code:} catalogue entry
\subsection{Definitions and mathematical statement}
A \emph{solution concept} for a partial differential equation specifies the admissible mathematical objects and the sense in which they satisfy the equation and any prescribed initial or boundary conditions.

\textbf{Source question.} What solution concepts are appropriate for the partial differential equations of fluid mechanics?

The supplied record poses a foundational research question rather than a quantified conjecture. It does not specify a particular equation or system, spatial dimension, spatial domain, time interval, initial or boundary conditions, regularity class, or criterion for appropriateness. Consequently, neither a definite class of admissible solutions nor a precise existence, uniqueness, stability, or selection assertion can be extracted without adding hypotheses absent from the source. These choices remain unresolved parts of the formulation.
\subsection{Short English statement}
The question asks which notions of solution are appropriate for fluid-mechanics PDEs. The equations, data, domains, and criteria for appropriateness are unspecified, so the supplied question does not determine a single precise mathematical assertion.
\subsection{Sources}
\begin{itemize}
\item Local registry: \texttt{data/unsolvedmath/problems.json}, permanent UnsolvedMath ID \(30006290\), source problem code \texttt{catalogue entry}; title: \emph{Solution Concepts for Fluid-Mechanics Equations}. The supplied record assigns the proposed year 2025 and names no individual proposer.
\item Supplied report citation: \emph{Mathematical Advances in Geophysical Fluid Dynamics} (2025), Oberwolfach Reports, DOI 10.4171/owr/2025/24, pp. 1215--1270. The original statement includes the page reference 1253. Supplied source URL and literature-source link, labeled ``original OWR report'': \url{https://doi.org/10.4171/owr/2025/24}.
\item Literature assessment dated exactly \texttt{2026-08-22}: the supplied curation labels the question open and describes it as a broad foundational question. It states that weak, dissipative, measure-valued, and selection concepts exist, but that no generally accepted complete solution concept is known. This is supplied dated curation derived from the OWR record, not an independently verified literature assessment; its time-sensitive status and claims require independent verification.
\item The supplied statement-status label is \texttt{corrected\_verified}. Its accompanying qualification says that only clear formula delimiters, notation, grammar, and typography were repaired, with mathematical content and scope preserved. This records the supplied editorial assessment, not independent mathematical verification.
\end{itemize}
\end{document}
